\begin{center}
	\zihao{4}\heiti
	%
	\LARGE{ 中国科学院大学\\ $\boldsymbol{2023}$年招收攻读硕士学位研究生入学统一考试试卷\\
	科目名称：电动力学（回忆版）}

\end{center}
\noindent
考生须知：\\
1.本试卷满分为150分,全部考试时间总计180分钟。\\
2.所有答案必须写在答题纸上,写在试题纸上或草稿纸上一律无效。

\vspace*{2ex}
\hrule
\vspace*{4ex}

\fancypagestyle{mypagestyle}{%
	\fancyhf{}% Clear header/footer
	\fancyfoot[OR]{第~\thepage~页，共2页}%
	\fancyfoot[EL]{第~\thepage~页，共2页}%
	\fancyfoot[OL]{科目名称：电动力学 }%
	\fancyfoot[ER]{科目名称：电动力学 }%
	\renewcommand{\headrulewidth}{0pt}% Header rule of .4pt
	\renewcommand{\footrulewidth}{0.5pt}% Header rule of .4pt
}
\fancypagestyle{plain}{%
	\fancyhf{}% Clear header/footer
	\fancyfoot[OR]{第~\thepage~页，共2页}%
	\fancyfoot[EL]{第~\thepage~页，共2页}%
	\fancyfoot[OL]{科目名称：电动力学 }%
	\fancyfoot[ER]{科目名称：电动力学 }%
}

\setcounter{page}{1}%设置页面计数器
\pagestyle{mypagestyle}

\onehalfspacing
% 排版结束，下面是题目
\vspace{0.4cm}

\noindent
一、选择题（共7题，每小题5分）
\vspace{0.2cm}

\noindent
1.\hspace{0.8em}试计算$\nabla \cdot [ (\Vec{a} \cdot \Vec{r}) \Vec{r} ]$，其中$\Vec{a}$为常矢量
\vspace{0.3cm}

A.\hspace{0.8em} 0 \hspace{2em} B.\hspace{0.8em} 2 $(\Vec{a} \cdot \Vec{r})$ \hspace{2em} C.\hspace{0.8em} 3 $(\Vec{a} \cdot \Vec{r})$ \hspace{2em} D.\hspace{0.8em} 4 $(\Vec{a} \cdot \Vec{r})$
\vspace{0.5cm}

\noindent
2.\hspace{0.8em}有两个电偶极子，$\Vec{P_1} = P_1 \Vec{e_z}$与$\Vec{P_2} = P_2 \Vec{e_z}$，则这两个电偶极子之间的静电力
\vspace{0.3cm}

A.\hspace{0.8em}为0 \hspace{8.9em}
B.\hspace{0.8em}始终为排斥力\vspace{0.2cm}

C.\hspace{0.8em}始终为吸引力\hspace{5em}
D.\hspace{0.8em}既可能为吸引力，又可能为排斥力\vspace{0.5cm}

\noindent
4.\hspace{0.8em}有一处电磁波入射到理想导体表面，则下列选项中错误的是\vspace{0.3cm}

A.\hspace{0.8em}电磁波无法射入理想导体，故$E = 0$，$B = 0$\vspace{0.2cm}

B.\hspace{0.8em}电磁波中的电场切向分量$E_{1t} = E_{2t} = 0$\vspace{0.2cm}

C.\hspace{0.8em}电磁波中的磁场切向分量$B_{1t} = B_{2t} = 0$\vspace{0.2cm}

D.\hspace{0.8em}电磁波中的磁场法向分量$B_{1n} = B_{2n} = 0$\vspace{0.5cm}

\noindent
7.\hspace{0.8em}两物体以$\displaystyle{ \frac{1}{2}C }$向相反的方向运动，一物体相对于另一物体的速度为\vspace{0.3cm}

A.\hspace{0.8em} 0 C \hspace{1em} B.\hspace{0.8em} 0.8 C \hspace{1em} C.\hspace{0.8em} 0.9 C \hspace{1em} D\hspace{0.8em} 0.95 C
\vspace{1cm}

\noindent
二、简答题（共20分）
\vspace{0.2cm}

\noindent
1.\hspace{0.8em}有一$K^{'}$系以速率$v$相对$K$系沿$x$轴正向运动，求二者之间的洛伦兹变化关系，并写出电动力学中的任意两个四维矢量。\vspace{0.5cm}

\noindent
2.\hspace{0.8em}试分析为什么在速度远小于光速c的情况下，带电粒子与电磁场作用的静磁力的影响远不如静电力。\vspace{0.5cm}

\noindent
3.\hspace{0.8em}写出介质中麦克斯韦方程组的积分形式，并写出各项同性介质中的
本构方程。\vspace{0.5cm}

\noindent
4.\hspace{0.8em}试用麦克斯韦方程推出电流连续方程$\nabla \cdot \Vec{J} + \displaystyle{\frac{\partial p}{\partial t}} = 0$，并根据能量守恒定律$\Vec{f} \cdot \Vec{v} + \nabla \cdot \Vec{S} + \displaystyle{\frac{\partial \omega}{\partial t}} = 0$，

\noindent
推导其中的能流密度与能量密度表达式。

\newpage

\noindent
三、\hspace{0.2em}（共25分）在均匀外磁场$\Vec{H}$中，有一磁导率为$\mu$的介质球
\vspace{0.2cm}

\noindent
1、用磁标势写出全空间的势，并写出磁感应强度；\vspace{0.1cm}

\noindent
2、写出介质球的磁矩。\vspace{4cm}


\noindent
四、\hspace{0.2em}（共25分）有一薄的导体球壳，其内径为$R_1$，外径为$R_2$。以球心为圆心，在$z$轴上$z=a$处放置一带电量为$Q$的电荷，导体球壳接地。
\vspace{0.2cm}

\noindent
1、试用镜像法，写出全空间的电势，并求出球壳内外表面上的面电荷密度。\vspace{0.1cm}

\noindent
2、设导体不接地，带有电荷量，求出全空间电势，并求出球壳内外表面上的面电荷密度。\vspace{4cm}

\noindent
四、\hspace{0.2em}（共25分）
有一无限长的矩形波导，在处插入一个板
\vspace{0.2cm}

\noindent
1、写出电磁波在矩形波导中满足的方程
\vspace{0.1cm}

\noindent
2、求电磁波在矩形波导中的边界条件
\vspace{0.1cm}

\noindent
3、电磁波在该波导中能以哪种形式的波模传播


\vspace{-2.5cm}
\begin{figure}[h!]
    \flushright
    \includegraphics[height=4cm]{test/2023/pic_1.png}
\end{figure}